\section{Partial data hiding}\label{sec:rotated-data-hiding}

The data-hiding game presented above was used to show that the
provers' measurement acts as identity on a subset of the provers'
qudits, and thus the prover learns no information from those
qudits. In particular, the measurement outcome of any $X$- or
$Z$-observable measurement on the qubits in the subset is hidden from
the prover. In this subsection, we generalize this idea to show how
to certify that certain \emph{partial} information about a register is hidden
from a prover. This test is a crucial component in our technique of
introspection, wherein two provers measure a shared EPR state to
sample from the joint distribution over questions of a classical
game. The partial data hiding test will prevent one prover from
learning the question sampled by the other prover.

\begin{notation}
  Given a set $v =\{v_1, \dots, v_k\}$ of $k$
  vectors in $\F_q^n$, denote their span by $V = \mathrm{span}(\{v_1,
  \dots, v_k\})$. The orthogonal complement of their span is the
  subspace $V^\perp = \{a: \forall i \in \{1, \dots, k\}, \langle a,
  v_i \rangle = 0 \}$. We denote by $\surfaces{v}$ the set of all
  affine subspaces parallel to $V$, i.e.\ sets of the form:
  \[ s = \{u + \lambda_1 v_1 + \dots + \lambda_k v_k: \lambda_1,
    \dots, \lambda_k \in \F_q\}.\]
  For a subspace $s \in
  \surfaces{v}$, the subspace projector $\Pi^{v}_s$ is the projector
  \[ \Pi^{v}_s = \sum_{w \in s} \proj{w}. \]
\end{notation}
\begin{lemma}
  \label{lem:subspace-x-z-commute}
  Given a set of vectors $\{v_1, \dots, v_k\}$, let
  \[ \tau^{X}_{[\forall i,  u \cdot v_i = a_i]} = \sum_{u: \forall i,
       u \cdot v_i = a_i} \tau^X_{u}. \]
  Then $\tau^{X}_{[\forall i,  u \cdot v_i = a_i]}$ commutes with
  $\Pi^{v}_s$ for all $s \in \surfaces{v}$.
\end{lemma}
\begin{proof}
  The proof is by calculation. 
  \begin{align*}
    \Pi^{v}_s \tau^{X}_{[\forall i,  u \cdot v_i = a_i]}
    &= \sum_{w \in s}\proj{w} \sum_{u: \forall i, u \cdot v_i = a_i}
      \tau^X_u \\
    &= \sum_{w \in s} \sum_{u: \forall i, u \cdot v_i = a_i} 
      \E_{\bb} \omega^{-\tr[\bb \cdot u]} \proj{w}  X(\bb). \\
    \intertext{We note an important fact: for any two
  outcomes $u, u'$ satisfying $u\cdot v_i = u'\cdot v_i = a_i$ for all
  $i$, the difference $u - u'$ must lie in $V^\perp$. Fixing some
    appropriate outcome vector $u_0$, we can then express the
    summation variable $u$ as $u_0 + x$ where $x$ runs over $V^\perp$: }
    &= \sum_{w \in s} \sum_{x \in V^\perp} \E_{\bb} \omega^{-\tr[\bb \cdot
      (u_0 + x)]} \proj{w} X(\bb) \\
    &=\sum_{w \in s} \sum_{x \in V^\perp} \E_{\bb} \omega^{-\tr[\bb \cdot
      (u_0 + x)]} \ket{w}\bra{w-\bb}.\\
    \intertext{Now, the summation over $x$ vanishes unless $\bb \in
    (V^\perp)^\perp = V$, by
    \Cref{fact:averages-to-zero-subspace}. This happens with
    probability $q^{k-n}$ which cancels out the factor of $q^{n-k}$
    from evaluating the sum over $x \in V^\perp$, yielding:}
    &= \sum_{w \in s} \E_{\bb \in V} \omega^{-\tr[\bb \cdot u_0]} \ket{w}
      \bra{w-\bb}. \\
    \intertext{Now, since $\bb \in V$, and the summation variable $w$ runs over an affine
    subspace parallel to $V$, we can shift it from $w$ to $w+\bb$, yielding}
    &= \sum_{w \in s} \E_{\bb \in V} \omega^{-\tr[\bb \cdot u_0]} \ket{w+\bb}
      \bra{w}. \\
    \intertext{Finally, we can perform the same manipulations in reverse:}
    &= \dots \\
    &= \tau^{X}_{[\forall i, u \cdot v_i = a_i]} \Pi^{v}_s. \qedhere
  \end{align*}
\end{proof}

% \begin{notation}
% \ainnote{DELETE THIS}
% Let $M$ be a matrix which operates on $\mathcal{H}_1 \otimes \mathcal{H}_{\mathrm{aux}}$.
% Define the notation
% \begin{equation*}
% \mathrm{hide}_{\mathcal{V}}(M)
% := \frac{1}{\tr(I_i)} \cdot I_i \otimes \tr_i(M).
% \end{equation*}
% \label{not:hide}
% \end{notation}


\begin{definition}\label{def:data-hide-def}
The \emph{partial data-hiding game} is given by~\Cref{fig:hide-observable}.
\end{definition}

{
\floatstyle{boxed} 
\restylefloat{figure}
\begin{figure}[htpb]
  Given a set $S$ of $k$-tuples of linearly independent set of vectors $v_1, \dots, v_k \in
  \F_q^{n}$ and a query string $x$. Sample $v = \{v_1, \dots, v_k\}$ uniformly from $S$.
  Flip an unbiased coin $\bb \sim \{0, 1\}$. Perform one of the following three tests with
  probability $1/3$ each.
  \begin{enumerate}
    \item Distribute the questions as follows:
      \begin{itemize}
      \item[$\circ$] Player~$\bb$: Give $(\bot, x,v)$; receive $(\varnothing, \ba_2)$.
      \item[$\circ$] Player~$\overline{\bb}$: give $(Z, x,v)$; receive
        $(\ba'_1, \ba'_2)$.
      \end{itemize}
      Accept if $\ba_2 = \ba_2'$.
    \item Distribute the questions as follows:
      \begin{itemize}
        \item[$\circ$] Player~$\bb$: Give $(\bot, x,v)$; receive $(\varnothing, \ba_2)$.
        \item[$\circ$] Player~$\overline{\bb}$: give $(\bot, x, \{X, v\})$; receive
        $(\varnothing, \ba'_2, \{\ba_{1,1}', \dots \ba_{1,k}'\})$. 
        \end{itemize}
        Accept if $\ba_2 = \ba_2'$.
      \item Distribute the questions as follows:
        \begin{itemize}
        \item[$\circ$] Player~$\bb$: Give $(X, \cdot)$; receive $(\ba_1, \cdot)$. (Here, ``$\cdot$" is the empty string.)
        \item[$\circ$] Player~$\overline{\bb}$: give $(\bot, \bot, \{X, v\})$; receive
        $(\varnothing,\varnothing, \{\ba_{1,1}', \dots \ba_{1,k}'\})$. 
        \end{itemize}
        Accept if $\ba'_{1,i} = v_i \cdot \ba_1$ for all $i \in \{1,
        \dots, k\}$.
      \item 
        Distribute the questions as follows:
        \begin{itemize}
          \item[$\circ$] Player~$\bb$: Give $(\bot, x, \{X, v\})$;
            receive $(\varnothing, \ba_2, \{\ba_{1,1}, \dots,
            \ba_{1,k}\})$.
          \item[$\circ$] Player~$\overline{\bb}$: give $(\bot,
            \bot, \{X, v\})$; receive $(\varnothing,
            \varnothing, \{\ba'_{1,1}, \dots, \ba'_{1,k}\})$.
          \end{itemize}
          Accept if $\ba_{1,i}  = \ba_{1,i}'$ for all $i \in \{1,
          \dots, k\}$.
      \end{enumerate}

\caption{The partial data-hiding game $\game_{\mathrm{hide}}(S, x)$.\label{fig:hide-observable}}
\end{figure}
}

  

\begin{theorem}  \label{thm:partial-data-hiding-game}
  Let $S$ be any set of $k$-tuples of vectors in
  $\F_q^n$, and let $x$ be an arbitrary query.
  \begin{itemize}
  \item[$\circ$]\textbf{Completeness:}
    Let $\calS_{\mathrm{partial}} = (\psi, M^{\bot, x,v})$ be a partial
    $(1,n,q)$-register strategy for~$\game_{\mathrm{hide}}(S,x)$,
    which is also a real commuting EPR strategy, and for which
    \[ M^{\bot, x,v}_{a_2} = \sum_{s \in \surfaces{v}} \Pi^{v}_s \ot A^{x,v,s}_{a_2}, \]
    for some measurement $A^{x,v,s}_{a_2}$ acting only on the
    $\reg{aux}$ register. Then
    there is a $(1,n,q)$-register strategy~$\calS$ extending
    $\calS_{\mathrm{partial}}$ 
    for which
    $\valstrat{\game_{\mathrm{hide}}}{\calS} = 1$.
    \item[$\circ$]\textbf{Soundness:}
      Let $\calS = (\psi, M)$ be a projective
      $(1,n,q)$-register strategy such that
      $\valstrat{\game_{\mathrm{hide}}(S,x)}{\calS}
      \geq 1 - \eps$. Then there exists an ideal measurement $M'^{\bot, x,
        v}_a$ with
      the property that
      \[ M'^{\bot, x,v }_a = \sum_{s \in \surfaces{v}} \Pi^{v}_s \ot
          M^{s,x, v}_a, \]
      such that  the measurement $M^{\bot, x,v}_a$ used by strategy $\calS$ in
      response to the query $x$  is close to $M'^{\bot, x,v}_a$:
      \[ (M_a^{\bot, x,v})_{\reg{Alice}} \otimes I_{\reg{Bob}} \approx_\eps (M'^{\bot, x,v}_a)_{\reg{Alice}} \otimes I_{\reg{Bob}}. \]
    \end{itemize}
\end{theorem}

To prove this theorem,  we will start with some basic facts about the subspace projector
  measurements. Let us denote the linear subspace spanned by the
  vectors $v_1, \dots, v_k$ by $V$.
  
  \begin{definition}
    For any distribution $\mathcal{U}$ over unitary matrices,
    the \emph{twirl by
      $\mathcal{U}$} is the linear operator $\twirl_{\mathcal{U}} :
    \calB((\mathbb{C}^q)^{\ot n}) \to \calB((\mathbb{C}^q)^{\ot n})$
    defined by
    \[ \twirl_{\mathcal{U}}(A) := \E_{\bU \sim \mathcal{U}} \left[ \bU A
        \bU^\dagger \right]. \]
  \end{definition}
  \begin{definition}
    Let $v = \{v_1, \dots, v_k\}$ be a set of linearly independent
    vectors over $\F_q$. Further let $\mathcal{V}$ be the uniform
    distribution over the set $\{X(a): a \in V \}$, $\mathcal{Z}$ be
    the uniform distribution over the set of all Pauli $Z$ operators
    $\{Z(a): a \in \F_q^n\}$, and $\mathcal{S}$ be the distribution
    over products $\bM\bN$ where $\bM$ is drawn from $\mathcal{V}$ and $\bN$
    from $\calZ$. Then the \emph{$v$-subspace twirl} is the
    twirl over $\calS$:
    \[ \twirl_{\calS} = \twirl_{\mathcal{V}} \circ \twirl_{\mathcal{Z}} \]
  \end{definition}
  \begin{proposition}
    \label{prop:hidden-subspace-twirl}
    Let $A$ be a Hermitian matrix and $v$ a set of $k$ vectors over $\F_q$. Then the $v$-subspace twirl of
    $A$ is a linear combination of projectors onto affine subspaces
    along $v$:
    \[ (\twirl_{\calS} \otimes \mathrm{id}_{\reg{aux}})(A) = \sum_{s \in \surfaces{v}}
      \Pi^{v}_{s} \ot (M_s)_{\reg{aux}} , \]
    for some choice of Hermitian matrices $M_s$ indexed by subspaces $s$.
  \end{proposition}
  \begin{proof}
    Start by decomposing $A$ into a linear combination of Pauli
    matrices:
    \[ A = \sum_{u,u'}  X(u) Z(u') \ot (A_{u,u'})_{\reg{aux}}. \]
    After the twirl over $\mathcal{Z}$, the only terms that survive
    are those with no $X$ part, i.e.
    \[ A' = (\twirl_{\mathcal{Z}} \otimes \mathrm{id}_{\mathrm{aux}})(A) = \sum_{u} Z(u) \ot (A_{0,u})_{\reg{aux}} \]
    Now if we perform the twirl over $\mathcal{V}$, we get
    \begin{align}
      (\twirl_{\mathcal{V}} \otimes \mathrm{id}_{\mathrm{aux}})(A')
      				&= \sum_u \E_{\ba \in V} X(\ba) Z(u) X(\ba)^\dagger \ot (A_{0,u})_{\reg{aux}}\nonumber\\
      				&= \sum_u  \E_{\ba \in V} \omega^{\tr[\langle
                                 \ba, u \rangle]} Z(u) \ot (A_{0,u})_{\reg{aux}}\nonumber \\
                               &= \sum_{u \in V^\perp}  Z(u) \ot (A_{0,u})_{\reg{aux}}\tag{\Cref{fact:averages-to-zero-subspace}}\nonumber\\
                               &= \sum_{u \in V^\perp}  \sum_{w} \omega^{\tr[\langle w, u\rangle]} \ket{w}\bra{w}  \ot (A_{0,u})_{\reg{aux}} \nonumber\\
                               &= \sum_{w}\ket{w}\bra{w}\ot \sum_{u \in V^\perp}  \omega^{\tr[\langle w, u\rangle]}    (A_{0,u})_{\reg{aux}}.\label{eq:let's-use-this-in-a-sec}
    \end{align}
    Now, consider a surface $s \in \surfaces{v}$.  For some~$x \in \F_q^n$,
    $s$ is the set of points written $w = x + v$, where $v \in V$.
    Then for any $u \in V^\perp$, $\langle w, u\rangle = \langle x + v, u \rangle = \langle x, u\rangle$,
    a quantity which depends only on the subspace and not on the point~$w$.
    Call this quantity $c_{s, u}$.
    As a result,
    \begin{equation*}
    \eqref{eq:let's-use-this-in-a-sec}
    = \sum_{s \in \surfaces{v}} \sum_{w \in s} \ket{w}\bra{w}\ot  \sum_{u \in V^\perp}c_{s, u} (A_{0,u})_{\reg{aux}}
    = \sum_{s \in \surfaces{v}} \Pi^{v}_{s} \ot (\hat{A}_s)_{\reg{aux}},
    \end{equation*}
    where $\hat{A}_s = \sum_{u \in V^\perp}c_{s, u} A_{0,u}$.
  \end{proof}
  
\begin{lemma}\label{lem:proj-to-obs-subspace}
Let $W \in \{X, Z\}$, and let $v = \{v_1, \dots, v_k\}$ be a set of
$k$ linearly independent vectors in $\F_q^n$ and $V$ be their span.
Suppose $\{M_{a_2}\}$ is a measurement for which
\begin{equation}\label{eq:m-commutes-subspace-proj}
(M_{a_2} \cdot (\tau^{W}_{[\forall i, v_i \cdot a_1 = a_{1,i}]} \otimes I_{\reg{aux}})) \otimes I_{\reg{Bob}}
\approx_{\delta} ((\tau^{W}_{[\forall i, v_i \cdot a_1 = a_{1,i}]} \otimes I_{\reg{aux}}) \cdot M_{a_2}) \otimes I_{\reg{Bob}},
\end{equation}
where
\[ \tau^{W}_{[\forall i, v_i \cdot a_1 = a_{1,i}]} = \sum_{a_1 :
    \forall i, v_i \cdot a_1 = a_{1,i}} \tau^W_{a_1}. \]
Then
\begin{equation*}
(M_{a_2} \cdot (W(u)\otimes I_{\reg{aux}})) \otimes I_{\reg{Bob}}
\approx_{\delta} ((W(u) \otimes I_{\reg{aux}}) \cdot M_{a_2}) \otimes I_{\reg{Bob}},
\end{equation*}
for a uniformly random $\bu$ drawn from $V$.
\end{lemma}
\begin{proof}
To start, given a set of outcomes $a_{1,1}, \dots, a_{1,k}$, suppose
$u$ and $u'$ are outcomes for a full $W$-basis measurement consistent with
these outcomes, i.e.\ $u$ and $u'$ are vectors such that for all $i$, $u \cdot
v_i = a_{1,i}$. Then it must hold that $u - u' \in V^\perp$. Using
this, the bound in~\Cref{eq:m-commutes-subspace-proj} becomes
\begin{equation}\label{eq:m-commutes-subspace-proj2}
  \sum_{a_2} \frac{1}{|V^\perp|} \sum_{u} \Big\| \sum_{w \in V^\perp} (M_{a_2} \cdot
  (\tau^W_{u+w} \otimes I_{\reg{aux}}) - (\tau^{W}_{u+w} \otimes I_{\reg{aux}}) \cdot M_{a_2}) \otimes I_{\reg{Bob}} \ket{\psi}\Big\|^2 \leq \delta,
\end{equation}
where the factor of $1/|V^\perp|$ is because each outcome $a_{1,1},
\cdots, a_{1,k}$ corresponds to $|V^\perp|$ different choices of $u$.
  
Our goal is to bound
\begin{equation}\label{eq:gonna-bound-this}
\E_{\bu\sim V} \sum_{a_2} \Vert (M_{a_2} \cdot (W(\bu)\otimes I_{\reg{aux}}) - (W(\bu)\otimes I_{\reg{aux}}) \cdot M_{a_2}) \otimes I_{\reg{Bob}} \ket{\psi} \Vert^2.
\end{equation}
by $\delta$.
To do so, for a fixed~$u$ we introduce the notation
\begin{align}
\Delta_{a_2}^{u}
&:= M_{a_2} \cdot (W(u) \otimes I_{\reg{aux}})- (W(u)\otimes I_{\reg{aux}}) \cdot M_{a_2}\\
&= \sum_{x \in \F_q^n} \omega^{\tr[u \cdot x]} (\underbrace{M_{a_2} \cdot (\tau^{W}_{x} \otimes I_{\reg{aux}})
	- (\tau^{W}_{x} \otimes I_{\reg{aux}}) \cdot M_{a_2}}_{\Delta_{a_2,x}}).\label{eq:pauli-obs-commutator}
\end{align}
We record the following identity, which follows
from~\Cref{eq:pauli-obs-commutator} and~\Cref{fact:averages-to-zero-subspace}:
\begin{equation*}
\E_{\bu \sim V } (\Delta_{a_2}^{\bu})^\dagger\cdot \Delta_{a_2}^{\bu}
= \E_{\bu \sim V} \sum_{x, x' \in \F_q^n} \omega^{\tr[\bu \cdot (x' - x)]} (\Delta_{a_2,x})^\dagger \Delta_{a_2,x'}
= \sum_{x \in \F_q^n} \sum_{w \in V^\perp} (\Delta_{a_2,x})^\dagger \Delta_{a_2,x+w}.
\end{equation*}
As a result,
\begin{align*}
\eqref{eq:gonna-bound-this}=
\E_{\bu} \sum_{a_2}\Vert (\Delta_{a_2}^{\bu} \otimes I_{\reg{Bob}}) \ket{\psi}\Vert^2
&= \E_{\bu} \sum_{a_2} \bra{\psi} (\Delta_{a_2}^{\bu})^\dagger \Delta_{a_2}^{\bu} \otimes I_{\reg{Bob}} \ket{\psi}\\
&=  \sum_{a_2} \sum_{x \in \F_q^n} \sum_{ w\in V^\perp} \bra{\psi}(\Delta_{a_2,x})^\dagger \Delta_{a_2,x+w} \otimes I_{\reg{Bob}} \ket{\psi}\\
&= \sum_{a_2, x} \frac{1}{|V^\perp|} \Big\| \sum_{w \in V^\perp} \Delta_{a_2,x+w} \otimes I_{\reg{Bob}} \ket{\psi} \Big\|^2,
\end{align*}
where the factor of $1/|V^\perp|$ is again to deal with overcounting. But this is at most $O(\delta)$, by~\Cref{eq:m-commutes-subspace-proj2}. This completes the proof.
\end{proof}

  
  \begin{lemma}\label{lem:commute-to-twirlU}
    Let $\{M_{a}\}$ be a measurement and $\calU$ be a distribution over
    unitaries,
    and suppose that for $U$ drawn uniformly from $\calU$,
    \begin{equation*}
      ((U^\dagger \otimes I_{\reg{aux}}) \cdot M_{a}) \otimes I_{\reg{Bob}}
      \approx_{\delta} (M_{a} \cdot (U^\dagger \otimes I_{\reg{aux}})) \otimes I_{\reg{Bob}},
    \end{equation*}
    where the distribution inherent in the $\approx_{\delta}$ notation
    is the uniform distribution over $\calU$.
    Then
    \begin{equation*}
      M_{a} \otimes I_{\reg{Bob}} \approx_{\delta} (\twirl_{\calU} \otimes I_{\reg{aux}})[M_{a}] \otimes I_{\reg{Bob}}.
    \end{equation*}
  \end{lemma}
  
  \begin{proof}
    By definition,
\begin{equation*}
(\twirl_1 \otimes I_{\reg{aux}})[M_{a}]
= \E_{\bU \sim \calU} [(\bU \otimes I_{\reg{aux}}) \cdot M_a \cdot (\bU^\dagger \otimes I_{\reg{aux}})].
\end{equation*}
Similarly,
\begin{equation*}
M_{a}
= \E_{\bU \sim \calU} [( \bU \otimes I_{\reg{aux}}) \cdot
(\bU^\dagger \otimes I_{\reg{aux}})  \cdot M_a].
\end{equation*}
As a result, if we set
\begin{equation*}
B(\bU)_a = (\bU^\dagger \otimes I_{\reg{aux}}) \cdot M_a - M_a \cdot (\bU^\dagger \ot I_{\reg{aux}}),
\end{equation*}
then
\begin{equation*}
\Delta_a:=
M_a - (\twirl_1 \otimes I_{\reg{aux}})[M_a]
= \E_{\bU \sim \calU} [\bU \cdot B(\bU)_a ].
\end{equation*}
We can therefore establish the lemma as follows:
\begin{align*}
\sum_a \Vert \Delta_a \otimes I_{\reg{Bob}} \ket{\psi} \Vert^2
& = \sum_a \Vert \E_{\bU \sim \calU} [ \bU \cdot B(\bU)_a ] \otimes I_{\reg{Bob}} \ket{\psi} \Vert^2\\
& \leq \E_{\bU} \sum_a \Vert (\bU \cdot B(\bU)_a ) \otimes I_{\reg{Bob}} \ket{\psi} \Vert^2\tag{Jensen's inequality}\\
& = \E_{\bU} \sum_a \Vert B(\bU)_a \otimes I_{\reg{Bob}} \ket{\psi} \Vert^2 \tag{$\bU$ is unitary}
\end{align*}
By assumption, this quantity is $O(\delta)$. This concludes the proof.
\end{proof}

\begin{proof}[Proof of \Cref{thm:partial-data-hiding-game}]
  We consider the completeness and soundness cases separately.
  \paragraph{Completeness}
  Let $\calS_{\mathrm{partial}} = (\psi, M^{\bot, x,v})$ be a partial
  $(k,n,q)$-register strategy for~$\game_{\mathrm{hide}}(S,x)$ which is also a real commuting EPR strategy, and for which
  the measurement $M^{\bot, x,v}_{a_2}$ has the form
  \[ M^{\bot, x,v}_{a_2} = \sum_{s \in \surfaces{v}} \Pi^{v}_s \ot A^{s,x,v}_{a_2} . \]
  To this strategy we will add matrices for the remaining questions.
  \begin{itemize}
  \item[$\circ$] Question $(Z, x)$: the measurement is
    \[ M^{(Z,x,v)}_{a_1, a_2} = \sum_{s \in \surfaces{v}} \Pi^{v}_s \cdot
      \tau^{Z}_{a_1} \ot A^{s,x,v}_{a_2}. \]
    This is a well-defined measurement as $\Pi^{v}_s$ is diagonal in the
    $Z$ basis and thus commutes with $\tau^Z_{a_1}$.
  \item[$\circ$] Question $(\bot, x, \{X, v_1, \dots, v_k\})$: the
    measurement is
    \[ M^{(\bot, x, \{X, v\})}_{\{a_{1,1}, \dots,
        a_{1,k}\}, a_2} = \sum_{s \in \surfaces{v}} \Pi^{v}_s \cdot
      \tau^{X}_{[\forall i, a_1 \cdot v_i = a_{1,i}]}  \ot
      A^{s,x,v}_{a_2}.\]
    This is a well-defined measurement as $\Pi^{v}_s$ commutes with
    $\tau^{X}_{[\forall i, a_1 \cdot v_i = a_{1,i}]}$
    by~\Cref{lem:subspace-x-z-commute}.
  \item[$\circ$] Question ($\bot, \bot, \{X, v_1, \dots, v_k\})$: the
    measurement is
    \[ M^{(\bot, \bot, \{X, v\})}_{\{a_{1,1}, \dots, a_{1,k},
        \varnothing} = \tau^X_{[\forall i, a_1 \cdot v_i = a_{1,i}]}
      \ot I. \]
  \item[$\circ$] Question $(X, \bot)$: the measurement is
    \[ M^{(X, \cdot)}_{a_1} = \tau^{X}_{a_1} \otimes I_{\reg{aux}}. \]
  \end{itemize}
This is a $(1,n,q)$-register strategy for~$\game_{\mathrm{hide}}$ by
design, and it is not hard to see that it achieves value $1$ on the
game. Assuming that the partial strategy $\cal{S}_{\mathrm{partial}}$
is a real commuting EPR strategy, it is not hard to see that the full
strategy above is also real (this is because if $M^{\bot, x,v}$ is real and
of the given form, then the matrices $A^{s,x,v}_{a_2}$ must also be
real). That the strategy is also commuting follows from the
description of $\game_{\mathrm{hide}}$. In particular, note that while
$M^{Z,x,v}_a$ does \emph{not} commute with $M^{\bot, x \{X, v\}}$ or
with $M^{X, \bot}$, the test never requires these measurements to be measured at the same time.

  \paragraph{Soundness}
   Recall that a strategy $\calS$ for this game
  consists of a state $\ket{\psi} = \ket{\epr^n_q} \ot \ket{\rmaux}$ and measurement operators of three types, corresponding to
  the four types of queries: $M^{(\bot,
    x,v)}_{\varnothing, a_2}$, $M^{(Z, x, v)}_{a_1, a_2}$, $M^{\bot, x, \{X, v\}, }_{\{a_{1,1}',
    \dots, a_{1,k}'\}, a_2'}$, and $M^{(X, \cdot)}_{a_1}$. 

  We start by analyzing the third and fourth parts of the test. The goal of these
  parts of the test is to certify that when given the query $(\{X, v\}, x)$, the prover returns $k$ answers $\ba_{1,1}',
  \dots, \ba_{1,k}'$ that are consistent with measuring the $X(v_1),
  \dots, X(v_k)$ observables on the state. We certify this in two
  stages. In part three of the test, we ask
  the first prover to do a complete measurement in the $X$ basis, and send
  the second prover the query $\{X, v\}$ indicating that it is to
  perform a partial $X$ measurement, and
  check consistency of outcomes. Importantly, in this part of the
  test, we \emph{cannot} send the second prover the query $x$, since
  the corresponding $\Pi^v_s$ measurement does not commute with the
  complete $X$ measurement performed by the first prover. Thus, in
  part four of the test, we send one prover the query $\{X, v\}$ and
  the other $(x, \{X, v\})$, and check consistency of their
  outcomes. 
  
  Since $\calS$ is a
  $(k,n,q)$-register strategy, \Cref{eq:in-math} implies that
  $M_{a_1}^{(X,\cdot)} = \tau^{X}_{a_1} \otimes I_{\reg{Bob}}$. We thus have
  from the third part of the test that
  \[ M^{\bot, \bot, \{X, v\}}_{\{a_{1,1}', \dots,
      a_{1,k}'\}} \ot I_{\reg{Bob}} \approx_\eps \tau^{X}_{[\forall i,
      u \cdot v_i = a'_{1,i}]} \ot I_{\reg{Bob}}.\]
  From the above equation and the fourth part of the test, we have
  \[ M^{\bot, x, \{X, v\}}_{\{a_{1,1}', \dots, a_{1,k}'\}} \ot
    I_{\reg{Bob}} \approx_\eps M^{\bot, \bot, \{X,v\}}_{\{a_{1,1}',
      \dots, a_{1,k}'\}} \ot I_{\reg{Bob}} \approx_\eps \tau^X_{[\forall i, u \cdot v_i =
      a'_{1,i}]} \ot I_{\reg{Bob}}. \]
  
  Next, we look at the first and second parts of $\game_{\mathrm{hide}}(S,x)$. These are
  essentially two instances of the commutation test. The first part of
  $\game_{\mathrm{hide}}$ certifies that the second outcome of $M^{Z,
    x,v}_{a_1, a_2}$ is consistent with $M^{\bot, x,v}_{\varnothing,
    a_2}$, and the hypothesis that the strategy $\calS$ is a
  $(1,n,q)$-register strategy tells us that the first outcome of
  $M^{Z,x,v}_{a_1, a_2}$ is consistent with $\tau^Z_{a_1} \ot I_{\reg{Bob}}$. Thus,
  applying the analysis of the commutation, it follows that
  \[ (M^{\bot, x,v }_{\varnothing, a_2} \cdot (\tau^Z_{a_1} \ot I_{\reg{aux}})) \otimes I_{\reg{Bob}} \approx_{\eps}
    ((\tau^Z_{a_1} \ot I_{\reg{Bob}})\cdot M^{\bot, x,v}_{\varnothing, a_2}) \otimes I_{\reg{Bob}}. \]
  A similar analysis for the second part of $\game_{\mathrm{hide}}(S)$
  certifies that
\[
((\tau^{X}_{[\forall i, a_1 \cdot v_i = a'_{1,i}]} \otimes
I_{\reg{aux}})\cdot M_{\varnothing, a_2}^{\bot,x, v}) \otimes I_{\reg{Bob}}
\approx_{\eps} (M_{\varnothing, a_2}^{\bot, x, v} \cdot(\tau^{X}_{[\forall i, a_1 \cdot v_i = a'_{1,i}]} \otimes I_{\reg{aux}})) \otimes I_{\reg{Bob}}.
\]

As a result, it holds that $W \in \{X,Z\}$, by \Cref{lem:proj-to-obs-subspace}, 
\begin{equation*}
(M_{\varnothing, a_2}^{\bot, x , v}\cdot (W(u)\otimes I_{\reg{aux}})) \otimes I_{\reg{Bob}}
\approx_{\eps} ((W(u)\otimes I_{\reg{aux}}) \cdot M_{\varnothing,
  a_2}^{\bot, x, v}) \otimes I_{\reg{Bob}}.
\end{equation*}
where if $W = X$, then $u$ is chosen uniformly over $V$, and
if $W = Z$, then $u$ is drawn uniformly from $\F_q^n$.
As a result, by~\Cref{fact:the-ol-pauli-swaperoonie} and \Cref{fact:add-a-proj},
\begin{align*}
(M_{\varnothing, a_2}^{\bot, x, v} \cdot (Z(u') X(u)\otimes I_{\reg{aux}})) \otimes I_{\reg{Bob}}
& \approx_{0}(M_{\varnothing, a_2}^{\bot,x,v} \cdot (Z(u')\otimes I_{\reg{aux}})) \otimes (X(-u)\otimes I_{\reg{aux}})\\
& \approx_{\eps}((Z(u')\otimes I_{\reg{aux}}) \cdot M_{\varnothing, a_2}^{\bot,x,v}) \otimes (X(-u)\otimes I_{\reg{aux}})\\
& \approx_{0}((Z(u')\otimes I_{\reg{aux}} )\cdot M_{\varnothing,a_2}^{\bot,x,v} \cdot (X(u) \otimes I_{\reg{aux}})) \otimes I_{\reg{Bob}}\\
& \approx_{\eps}((Z(u') X(u)\otimes I_{\reg{aux}}) \cdot
                                                                                                                                    M_{\varnothing,
                                                                                                                                    a_2}^{\bot,x,
                                                                                                                                    v})
                                                                                                                                    \otimes I_{\reg{Bob}},
\end{align*}
on the distribution where $v$ is chosen from $S$, and then $\bu \sim V$, $\bu' \sim \F_q^n$.
Applying \Cref{lem:commute-to-twirlU} and \Cref{prop:hidden-subspace-twirl},
we can therefore conclude that
\begin{equation*}
M_{\varnothing, a_2}^{\bot, x,v} \otimes I_{\reg{Bob}}
\approx_{\eps} (\twirl_{\calS} \otimes I_{\reg{aux}})[M_{\varnothing, a_2}^x]
\otimes I_{\reg{Bob}}
= \left(\sum_{s \in \surfaces{v}} \Pi^{v}_s \ot (A^{x,v,s}_{a_2})_{\reg{aux}}\right)
\otimes I_{\reg{Bob}}, 
\end{equation*}
for some choice of measurements $A^{x,v,s}_{a_2}$ on the aux register.
\end{proof}

%%% Local Variables:
%%% mode: latex
%%% TeX-master: "main"
%%% End:
%%%---------- close: partial-data-hiding

\part{$\neexp$ protocol}

\label{part:neexp}

%%%%%%%%%%% jump to little-nexp-protocol
%%%---------- open: little-nexp-protocol
%!TEX root = main.tex

